% GENERATED FILE. Edit canonical lessons, then run npm run generate:books.
% canonical-source-hash: fnv1a64:7a746be0843d4a62
% canonical-lessons: JA-C495-nioi, JA-C495-kaori, JA-C495-fuumi, JA-C495-shinsen, JA-C495-assari

\chapter{Smell and Flavour}
\label{ch:ja-smell-and-flavour}
\hlchaptermodality{495}

\begin{chapteropening}
\textbf{By the end of this chapter:} \emph{I can describe how food smells and tastes.}

Complete the last lesson of chapter 495: I can describe how food smells and tastes.
\end{chapteropening}

\section[\texorpdfstring{\ja{におい} (nioi)}{におい (nioi)}]{\ja{におい} (nioi) — a smell}

\label{lesson:JA-C495-nioi}



\begin{quote}
Before the new one: say the Japanese for a bamboo shoot, then the Japanese for a shiitake mushroom.
\end{quote}

\subsection*{You'll want to know: \ja{におい}}
\textbf{\ja{におい}} — \emph{nioi} — \textquotedblleft{}a smell\textquotedblright{}.

\begin{quote}
\textbf{\ja{だいどころから} \ja{いい} \ja{においが} \ja{します。}} — \emph{daidokoro kara ii nioi ga shimasu} — \textquotedblleft{}A good smell is coming from the kitchen.\textquotedblright{}
\end{quote}

\begin{cousinweb}
\textbf{\ja{におい}} (\emph{nioi}) is a native Japanese word (wago), inherited rather than borrowed from Chinese; its kanji was attached later as a kun'yomi, and its older origin is not settled (JMdict, KANJIDIC2).
\end{cousinweb}

\begin{grammarlens}[title={What you've built}]
One more word for describing what you lived and felt.
\end{grammarlens}

\subsection*{Your turn}
\begin{itemize}
\raggedright
  \item \emph{Say it:} \emph{nioi}
  \item \emph{Say it:} \emph{nioi}, once more
  \item \emph{Say it:} \emph{daidokoro kara ii nioi ga shimasu} — A good smell is coming from the kitchen.
  \item \emph{Recall:} say the Japanese for a bamboo shoot, then the Japanese for a shiitake mushroom, then say \emph{nioi} again
\end{itemize}

\begin{tcolorbox}[breakable,colback=teal!4,colframe=teal!35!black,title={Before you move on}]
What does \ja{におい} mean? (\textquotedblleft{}A smell\textquotedblright{}.) Say it once more.
\end{tcolorbox}
\section[\texorpdfstring{\ja{かおり} (kaori)}{かおり (kaori)}]{\ja{かおり} (kaori) — a fragrance; an aroma}

\label{lesson:JA-C495-kaori}



\begin{quote}
Before the new one: say the Japanese for a shiitake mushroom, then the Japanese for a smell.
\end{quote}

\subsection*{You'll want to know: \ja{かおり}}
\textbf{\ja{かおり}} — \emph{kaori} — \textquotedblleft{}a fragrance; an aroma\textquotedblright{}.

\begin{quote}
\textbf{\ja{この} \ja{おちゃは} \ja{かおりが} \ja{いいです。}} — \emph{kono ocha wa kaori ga ii desu} — \textquotedblleft{}This tea has a lovely aroma.\textquotedblright{}
\end{quote}

\begin{cousinweb}
\textbf{\ja{かおり}} (\emph{kaori}) is a native Japanese word (wago), inherited rather than borrowed from Chinese; its kanji was attached later as a kun'yomi, and its older origin is not settled (JMdict, KANJIDIC2).
\end{cousinweb}

\begin{grammarlens}[title={What you've built}]
One more word for describing what you lived and felt.
\end{grammarlens}

\subsection*{Your turn}
\begin{itemize}
\raggedright
  \item \emph{Say it:} \emph{kaori}
  \item \emph{Say it:} \emph{kaori}, once more
  \item \emph{Say it:} \emph{kono ocha wa kaori ga ii desu} — This tea has a lovely aroma.
  \item \emph{Recall:} say the Japanese for a shiitake mushroom, then the Japanese for a smell, then say \emph{kaori} again
\end{itemize}

\begin{tcolorbox}[breakable,colback=teal!4,colframe=teal!35!black,title={Before you move on}]
What does \ja{かおり} mean? (\textquotedblleft{}A fragrance; an aroma\textquotedblright{}.) Say it once more.
\end{tcolorbox}
\section[\texorpdfstring{\ja{ふうみ} (fūmi)}{ふうみ (fūmi)}]{\ja{ふうみ} (fūmi) — flavour; taste}

\label{lesson:JA-C495-fuumi}



\begin{quote}
Before the new one: say the Japanese for a smell, then the Japanese for a fragrance.
\end{quote}

\subsection*{You'll want to know: \ja{ふうみ}}
\textbf{\ja{ふうみ}} — \emph{fūmi} — \textquotedblleft{}flavour; taste\textquotedblright{}.

\begin{quote}
\textbf{\ja{この} \ja{みそしるは} \ja{ふうみが} \ja{いいです。}} — \emph{kono misoshiru wa fūmi ga ii desu} — \textquotedblleft{}This miso soup has a good flavour.\textquotedblright{}
\end{quote}

\begin{cousinweb}
\textbf{\ja{ふうみ}} (\emph{fūmi}) is Sino-Japanese (kango): two Chinese characters read with their on'yomi, readings borrowed from Chinese. They are \emph{fū}, \textquotedblleft{}wind\textquotedblright{} (Middle Chinese \emph{pjuwng}) and \emph{mi}, \textquotedblleft{}flavor\textquotedblright{} (\emph{mj+jH}) (KANJIDIC2; Middle Chinese in Baxter's transcription).
\end{cousinweb}

\begin{grammarlens}[title={What you've built}]
One more word for describing what you lived and felt.
\end{grammarlens}

\subsection*{Your turn}
\begin{itemize}
\raggedright
  \item \emph{Say it:} \emph{fūmi}
  \item \emph{Say it:} \emph{fūmi}, once more
  \item \emph{Say it:} \emph{kono misoshiru wa fūmi ga ii desu} — This miso soup has a good flavour.
  \item \emph{Recall:} say the Japanese for a smell, then the Japanese for a fragrance, then say \emph{fūmi} again
\end{itemize}

\begin{tcolorbox}[breakable,colback=teal!4,colframe=teal!35!black,title={Before you move on}]
What does \ja{ふうみ} mean? (\textquotedblleft{}Flavour; taste\textquotedblright{}.) Say it once more.
\end{tcolorbox}
\section[\texorpdfstring{\ja{しんせん} (shinsen)}{しんせん (shinsen)}]{\ja{しんせん} (shinsen) — fresh}

\label{lesson:JA-C495-shinsen}



\begin{quote}
Before the new one: say the Japanese for a fragrance, then the Japanese for flavour.
\end{quote}

\subsection*{You'll want to know: \ja{しんせん}}
\textbf{\ja{しんせん}} — \emph{shinsen} — \textquotedblleft{}fresh\textquotedblright{}.

\begin{quote}
\textbf{\ja{いちばの} \ja{さかなは} \ja{しんせんです。}} — \emph{ichiba no sakana wa shinsen desu} — \textquotedblleft{}The fish at the market is fresh.\textquotedblright{}
\end{quote}

A na-adjective: \ja{な} before a noun.

\begin{cousinweb}
\textbf{\ja{しんせん}} (\emph{shinsen}) is Sino-Japanese (kango): two Chinese characters read with their on'yomi, readings borrowed from Chinese. They are \emph{shin}, \textquotedblleft{}new\textquotedblright{} (Middle Chinese \emph{sin}) and \emph{sen}, \textquotedblleft{}fresh\textquotedblright{} (\emph{sjen}) (KANJIDIC2; Middle Chinese in Baxter's transcription).
\end{cousinweb}

\begin{grammarlens}[title={What you've built}]
One more word for describing what you lived and felt.
\end{grammarlens}

\subsection*{Your turn}
\begin{itemize}
\raggedright
  \item \emph{Say it:} \emph{shinsen}
  \item \emph{Say it:} \emph{shinsen}, once more
  \item \emph{Say it:} \emph{ichiba no sakana wa shinsen desu} — The fish at the market is fresh.
  \item \emph{Recall:} say the Japanese for a fragrance, then the Japanese for flavour, then say \emph{shinsen} again
\end{itemize}

\begin{tcolorbox}[breakable,colback=teal!4,colframe=teal!35!black,title={Before you move on}]
What does \ja{しんせん} mean? (\textquotedblleft{}Fresh\textquotedblright{}.) Say it once more.
\end{tcolorbox}
\section[\texorpdfstring{\ja{あっさり} (assari)}{あっさり (assari)}]{\ja{あっさり} (assari) — light (of food); plain and simple}

\label{lesson:JA-C495-assari}



\begin{quote}
Before the new one: say the Japanese for flavour, then the Japanese for fresh.
\end{quote}

\subsection*{You'll want to know: \ja{あっさり}}
\textbf{\ja{あっさり}} — \emph{assari} — \textquotedblleft{}light (of food); plain and simple\textquotedblright{}.

\begin{quote}
\textbf{\ja{この} \ja{りょうりは} \ja{あっさりしています。}} — \emph{kono ryōri wa assari shiteimasu} — \textquotedblleft{}This dish is light and simple.\textquotedblright{}
\end{quote}

\begin{cousinweb}
\textbf{\ja{あっさり}} (\emph{assari}) is a native mimetic word (gitaigo), imitating the impression it describes rather than borrowed from Chinese; its older origin is not settled (JMdict).
\end{cousinweb}

\begin{grammarlens}[title={What you've built}]
One more word for describing what you lived and felt.
\end{grammarlens}

\subsection*{Your turn}
\begin{itemize}
\raggedright
  \item \emph{Say it:} \emph{assari}
  \item \emph{Say it:} \emph{assari}, once more
  \item \emph{Say it:} \emph{kono ryōri wa assari shiteimasu} — This dish is light and simple.
  \item \emph{Recall:} say the Japanese for flavour, then the Japanese for fresh, then say \emph{assari} again
\end{itemize}

\begin{tcolorbox}[breakable,colback=teal!4,colframe=teal!35!black,title={Before you move on}]
What does \ja{あっさり} mean? (\textquotedblleft{}Light (of food); plain and simple\textquotedblright{}.) Say it once more.
\end{tcolorbox}
